\section{Oscillatory integrals and Bochner--Riesz multipliers}
\label{sec:multiplier}

We now turn packet growth into the multiplier bound. The first step is
an oscillatory integral estimate with loss $\lambda^\nu$ for every
$\nu>0$, uniform over fixed phase and amplitude classes. After a
change of variables, this gives norm
$O_{\delta,\nu}(\lambda^{-\delta+\nu})$ for a Bochner--Riesz kernel
shell of radius $\lambda$. Choosing $\nu<\delta$ makes the dyadic
shell sum converge. We keep the data bounds explicit because the
shell amplitudes vary with the scale.

\begin{theorem}[Oscillatory integral estimate]\label{thm:oscillatory-estimate}
Fix bounded input and output regions $K_v,K_x\subset\RR^2$.
For the input patch $K_v$, fix a class of smooth real phases satisfying
the phase assumptions of Section~\ref{sec:packet-structure}, with
common domain $[0,1]\times\RR^2\times U_v$, where $U_v$ is an open
neighborhood of $K_v$.  All enlarged patches required by the
decomposition, expressed in the original $v$ coordinates, are contained
in $U_v$ and obey common bounds.  The translated systems are considered
on the corresponding bounded rescaled patches described in that section,
with the resulting uniform bounds.  Throughout the phase domain,
\[
 \partial_v\Phi(c,x,v)=\zeta(v,x-cv).
\]
For each phase and each required original input patch $V$, the matrix
$\partial_p\zeta$ is symmetric and satisfies either
\[
 a_0I\le\partial_p\zeta(v,p)\le a_1I
 \qquad(v\in V,\ p\in\RR^2),
\]
or the same inequality with $-\partial_p\zeta$ in place of
$\partial_p\zeta$, with common constants $a_0>0$ and $a_1$ for the
class.  On these original patches, every mixed $v,p$ derivative of total
order $j$ of this matrix is bounded by $A^{j+1}(j!)^3$ for a common
$A\ge1$.  Fix these constants and all the seminorm bounds of a class of
smooth amplitudes $B$ supported in
$[0,1]\times K_x\times K_v$.  For every $\nu>0$ there is a constant
$C_\nu$, depending only on these fixed data and $\nu$, such that, for
every $\lambda\ge2$, every phase and amplitude in the specified classes,
and every $f\in L^3(\RR^2)$ supported in $K_v$,
\begin{equation}\label{eq:oscillatory-estimate}
 \left\|E_\lambda f\right\|_{L^3([0,1]\times K_x)}
 \le C_\nu\lambda^{-1+\nu}\|f\|_3,
 \qquad
 E_\lambda f(c,x)=\int e^{i\lambda\Phi(c,x,v)}B(c,x,v)f(v)\,dv.
\end{equation}
\end{theorem}

\begin{proof}
Fix $0<\gamma<1$ and a loss $\eta>0$.  Apply
Theorem~\ref{thm:packet-growth} with
\[
 \eps=\lambda^{-1},\qquad \Planck=t=L=1.
\]
Thus its finite conclusion is $P_1\le K_0+\gamma/2+\eta$ for all
sufficiently large $\lambda$, uniformly in the data under consideration.
Here and until the endpoint sum below, logarithms in scores have base
$\eps^{-1}$.

Suppose first that $|f|\le\ind_F$, where $F\subset K_v$ is measurable.
The case $f=0$ is immediate.  At the root the packet radius is
$\rho_0=\lambda^{-1/2}$.  Write $m_O$ for the canonical root masses.
The Bessel bound of Lemma~\ref{lem:packet-bessel} gives
\[
 M:=\sum_Om_O\le C|F|.
\]
If $D$ is any admissible root aperture of widths $b\le a\le1/2$, its
velocity centers lie in a rectangle with side lengths comparable to
$\eps^b,\eps^a$.  All the corresponding test windows lie in the
$100\rho_0$ enlargement of this rectangle.  Since
$\rho_0\le\eps^a\le\eps^b$, that enlargement has area at most
$C\eps^{a+b}$.  Applying the same Bessel estimate to $f$ restricted to
the enlarged rectangle yields
\begin{equation}\label{eq:restricted-root-cap}
 M_D:=\sum_{O\text{ compatible with }D}m_O
 \le C\eps^{a+b}.
\end{equation}
This estimate holds separately for each label, including labels that
are assigned stochastically.  Its constant includes the fixed lattice
and window multiplicities.

For any law of $(O,D)$ supported on compatible pairs, log-sum gives
\[
 \Ent(O)+\EE\logeps{m_O}\le\logeps M,
 \qquad
 \Ent(O\mid D)+\EE\logeps{m_O}\le\EE\logeps{M_D}.
\]
Consequently every root score is at most
\[
 \logeps{|F|}
 +\tfrac12\EE\bigl[w(b,a)-a-b\bigr]
 +\frac{C}{\log\lambda}
 \le \logeps{|F|}+\frac{C}{\log\lambda},
\]
because $w(b,a)-a-b=-\gamma(a-b)\le0$.  The norm option in $K_0$
has the same upper bound: indeed
\[
 \|f\|_2^3\le |F|^{3/2}\le |K_v|^{1/2}|F|.
\]
We have therefore proved
\begin{equation}\label{eq:restricted-root-score}
 K_0\le\logeps{|F|}+\frac{C}{\log\lambda}.
\end{equation}

At $t=\Planck=1$ the velocity radius is one.  Cover the input region
by a fixed number of such velocity patches, and partition the output
region into boxes with time and position side lengths comparable to
$\lambda^{-1}$.  We may work with one input patch at a time.  For each
output box let $m_O$ be a test mass controlling the supremum of
$|E_\lambda f|^2$ on that box, up to a fixed factor.  These tests are
admissible endpoint masses.  To see the asserted uniformity, if
$(C,X)$ is the box anchor, absorb
\[
 \exp\!\left(i\lambda\bigl[
 \Phi(c,x,v)-\Phi(C,X,v)
 -\Phi(c,x,v_*)+\Phi(C,X,v_*)\bigr]\right)
\]
into the amplitude, where $v_*$ is fixed in the input patch.  Every
positive-order $v$ derivative of the phase difference is
$O(\lambda^{-1})$, by the defining identity for $\partial_v\Phi$ and
the uniform bounds on $\partial_p\zeta$.  The displayed factor thus
has uniformly bounded derivatives.  Sitewise choices approaching a
supremum are permitted by the packet theorem; no regularity between
the choices at distinct boxes is required.

Only the aperture widths $b=a=0$ are allowed at the endpoint.  Choose
their labels with a fixed orientation.  For any endpoint law its score
is at least $\Ent(O)+\tfrac32\EE\logeps{m_O}$.  Taking the law
proportional to $m_O^{3/2}$, or omitting zero masses, gives the exact
log-sum value
\[
 \Ent(O)+\tfrac32\EE\logeps{m_O}
 =\logeps{\sum_Om_O^{3/2}}.
\]
The packet theorem and \eqref{eq:restricted-root-score} imply
\begin{equation}\label{eq:endpoint-packet-sum}
 \sum_Om_O^{3/2}\le C_{\gamma,\eta}
 \lambda^{\gamma/2+\eta}|F|.
\end{equation}
Each output box has volume $O(\lambda^{-3})$.  Summing the integrals
over boxes, and then the fixed number of input patches, proves
\begin{equation}\label{eq:restricted-oscillatory}
 \|E_\lambda f\|_3
 \le C_{\gamma,\eta}\lambda^{-1+\gamma/6+\eta/3}|F|^{1/3}
 \qquad (|f|\le\ind_F).
\end{equation}
All constants are uniform over the fixed classes.  In particular, the
passage from limiting packet growth to the finite bound used here is
uniform: otherwise a sequence of offending scales, inputs, phases,
amplitudes, and laws would contradict the sequence quantifier in
Theorem~\ref{thm:packet-growth}.

For a general input normalize $\|f\|_3=1$.  The two tails satisfy
\[
 \|f\ind_{\{|f|\le\lambda^{-10}\}}\|_1
 \le |K_v|\lambda^{-10},\qquad
 \|f\ind_{\{|f|>\lambda^{10}\}}\|_1\le\lambda^{-20}.
\]
The trivial $L^1$ to bounded-region $L^3$ estimate controls their
contributions.  Split the remaining input into $N=O(\log\lambda)$
dyadic levels $f_k$, with $2^k<|f|\le2^{k+1}$ on $F_k$.  Equation
\eqref{eq:restricted-oscillatory}, followed by H\"older in $k$, gives
\[
 \left\|\sum_k E_\lambda f_k\right\|_3
 \le C_{\gamma,\eta}\lambda^{-1+\gamma/6+\eta/3}
 \sum_k2^{k+1}|F_k|^{1/3}
 \le 2C_{\gamma,\eta}\lambda^{-1+\gamma/6+\eta/3}N^{2/3}.
\]
Choose $\gamma,\eta>0$ so that $\gamma/6+\eta/3<\nu/2$ and absorb
$N^{2/3}$ in the remaining positive power of $\lambda$.  Bounded values
of $\lambda$ are covered by the trivial estimate.  Homogeneity proves
\eqref{eq:oscillatory-estimate} for every input.
\end{proof}

\subsection{The radial kernel and its dyadic pieces}

Use the Fourier convention
$\widehat f(\xi)=\int e^{-2\pi i x\cdot\xi}f(x)\,dx$.  For real
$\delta>0$ put
\[
 m_\delta(\xi)=(1-|\xi|^2)_+^\delta,
 \qquad K_\delta=\mathcal F^{-1}m_\delta.
\]

\begin{lemma}[Endpoint expansion of the kernel]\label{lem:br-kernel}
The kernel is bounded on $\RR^3$.  For $R\ge1$ it has an expansion
\begin{equation}\label{eq:br-kernel-symbol}
 K_\delta(Re_3)=e^{2\pi iR}b_+(R)+e^{-2\pi iR}b_-(R),
 \qquad
 |b_\pm^{(j)}(R)|\le C_{\delta,j}R^{-2-\delta-j}
 \quad(j\ge0).
\end{equation}
In particular $K_\delta\in L^1(\RR^3)$ whenever $\delta>1$.
\end{lemma}

\begin{proof}
Boundedness follows from $m_\delta\in L^1$.  Integrating first over
the disks perpendicular to $e_3$ gives the exact formula
\[
 K_\delta(Re_3)=\frac{\pi}{\delta+1}
 \int_{-1}^1e^{2\pi iRz}(1-z^2)^{1+\delta}\,dz.
\]
Choose a smooth partition into neighborhoods of the two endpoints and
a compact subinterval of $(-1,1)$.  At $z=1$ write $t=1-z$ and
extract $e^{2\pi iR}$.  Its remaining integral has amplitude
$t^{1+\delta}a(t)$, with $a$ smooth and supported in a fixed small
interval $[0,t_0)$.  Its $j$th $R$ derivative has amplitude a constant
times $t^{1+\delta+j}a(t)$.  The part $t\le R^{-1}$ is bounded by
$CR^{-2-\delta-j}$.  On a smooth dyadic shell $t\asymp h\ge R^{-1}$,
$N$ integrations by parts give the bound
\[
 C_{\delta,j,N}R^{-N}h^{2+\delta+j-N}.
\]
For an integer $N>2+\delta+j$ the sum over these shells is at most
$C_{\delta,j}R^{-2-\delta-j}$.  The same argument at $z=-1$ gives
the other symbol.  The middle integral and every derivative decrease
faster than any inverse power of $R$; after multiplication by
$e^{-2\pi iR}$ it can be included in $b_+$.  This proves
\eqref{eq:br-kernel-symbol}.  Radial integration of its absolute-value
bound gives $\int_1^\infty R^2R^{-2-\delta}\,dR<\infty$ when
$\delta>1$.
\end{proof}

Take a smooth radial partition of unity into a bounded region and
shells $|x|\asymp\lambda$, $\lambda=2^j\ge2$.  The bounded piece
of $K_\delta$ is integrable.  By Lemma~\ref{lem:br-kernel}, each shell
is the sum of two kernels
\[
 \lambda^{-2-\delta}e^{\pm2\pi i|x|}
 a_{\lambda,\pm}(x/\lambda),
\]
where the amplitudes are supported in a fixed annulus and all their
derivatives are bounded uniformly in $\lambda$.  In the convolution
operator set $x=\lambda z$, $y=\lambda y'$.  The factor from volume
is $\lambda^3$, so the rescaled operator is
\begin{equation}\label{eq:br-rescaled-shell}
 \lambda^{1-\delta}A_{\lambda,\pm}f(z),\qquad
 A_{\lambda,\pm}f(z)=
 \int e^{\pm2\pi i\lambda|z-y|}a_{\lambda,\pm}(z-y)f(y)\,dy.
\end{equation}
The input and output dilation factors cancel in its $L^3$ operator
norm.  We next prove
\begin{equation}\label{eq:br-unit-shell}
 \|A_{\lambda,\pm}\|_{L^3(\RR^3)\to L^3(\RR^3)}
 \le C_\nu\lambda^{-1+\nu} \qquad(\nu>0).
\end{equation}

\subsection{The distance phase}

Partition input and output space into unit cubes with smooth cutoffs
of bounded overlap.  The annular support in \eqref{eq:br-rescaled-shell}
allows only a bounded number of input neighbors for each output cube,
and conversely.  For one such pair, translate both cubes together to
a fixed bounded region.  A finite smooth angular partition and fixed
rotations reduce the phase to regions where
\[
 0<d_0\le z_3-y_3\le d_1<\infty,
\]
with fixed constants. The dyadic shell operators have the
Carleson--Sj\"olin distance-phase form treated by
\citet{GuoOhWangWuZhang2025}; we use their subsequent
pseudoconformal change of variables. Fix the slice variable $y_3$, and set
\[
 s=(z_3-y_3)^{-1},\qquad Z=s(z_1,z_2),\qquad v=(y_1,y_2).
\]
Then, for $g(q)=\sqrt{1+|q|^2}$,
\[
 |z-y|=s^{-1}g(Z-sv),\qquad
 \left|\det\frac{\partial(Z,s)}{\partial z}\right|=s^4.
\]
Choose a fixed positive interval $[s_*,s_*+D]$ containing the range
of $s$, and put $s=s_*+Dc$, $Z=Dx$.  Both the coordinate change and
its inverse have bounded derivatives and Jacobians on the relevant
regions, uniformly in the slice.  Direct differentiation gives
\begin{align}
 \Phi(c,x,v)&=\pm2\pi(s_*+Dc)^{-1}
 g\bigl(Dx-(s_*+Dc)v\bigr),\label{eq:distance-phase}\\
 \partial_v\Phi(c,x,v)&=\zeta(v,x-cv),\qquad
 \zeta(v,p)=\mp2\pi\nabla g(Dp-s_*v),\label{eq:distance-zeta}\\
 \partial_p\zeta(v,p)&=\mp2\pi D\nabla^2g(Dp-s_*v).
 \label{eq:distance-hessian}
\end{align}
The eigenvalues of $\nabla^2g(q)$ are
$(1+|q|^2)^{-1/2}$ and $(1+|q|^2)^{-3/2}$.  They are bounded away
from zero on the fixed argument region of the integral, but the smaller
eigenvalue tends to zero at infinity.  The packet decomposition uses
global definiteness when solving \eqref{eq:packet-gradient-matching}
for every Fourier mode.  We therefore modify $g$ outside the argument
region to obtain the global hypothesis of
Theorem~\ref{thm:oscillatory-estimate}.

Choose $R_0>1$ so that all arguments in \eqref{eq:distance-phase} have
norm less than $R_0$.  Let $\chi$ be a nondecreasing smooth function,
zero on $(-\infty,R_0]$, one on $[R_0+1,\infty)$, with
\[
 |\chi^{(j)}(r)|\le C^{j+1}(j!)^2.
\]
Such a function is obtained by integrating and normalizing the bump
$\exp[-1/t-1/(1-t)]$ on $0<t<1$, extended by zero.  For completeness,
Cauchy's estimate on a complex disk of radius a small fixed multiple
of $t$ gives a derivative bound
$j!(C/t)^j e^{-c/t}$ near $t=0$.  Maximizing $t^{-j}e^{-c/t}$ gives
$C^{j+1}(j!)^2$; the estimate at $t=1$ is identical, and away from
the endpoints analyticity gives the bound directly.

Let $h(r)=0$ for $r\le R_0$ and, for $r>R_0$, let
\[
 h(r)=\int_{R_0}^r(r-u)\chi(u)\,du.
\]
Replace $g(q)$ by $\widetilde g(q)=g(q)+h(|q|)$.  The added Hessian
has radial and tangential eigenvalues $\chi(r)$ and $h'(r)/r$, both
nonnegative and at most one.  For $r\ge2(R_0+1)$ they are at least
$1/2$.  On the remaining compact region the original Hessian is
uniformly positive.  Thus
\[
 c_0 I\le\nabla^2\widetilde g\le C_0 I\quad\hbox{on }\RR^2.
\]
The modification vanishes on the argument region.  Its Hessian has
derivative bounds $A_0^{j+1}(j!)^3$: on the transition annulus this
follows by the chain rule from the displayed Gevrey bounds and the
analyticity of $q\mapsto|q|$ away from zero.  More explicitly, grouping
the terms in a derivative of order $j$ by the $k$ derivatives landing
on the outer scalar function bounds their sum by
\[
 C^j j!\sum_{k=1}^j C^k k!\binom{j-1}{k-1}
 \le C_1^{j+1}(j!)^2.
\]
The analytic factors $q/|q|$ in the radial Hessian obey the same bound
after the product rule.  Outside the transition annulus,
$h(r)$ is a quadratic polynomial in $r$ plus a linear term and a
constant, whose radial Hessian has the same bounds.  Near the origin
there is no modification.  The original $g$ has uniform analytic
Hessian bounds.  Formula \eqref{eq:distance-hessian}, with
$\widetilde g$ in place of $g$, therefore satisfies all the required
global hypotheses, including the mixed $v,p$ derivative bounds.

For each fixed $y_3$, Theorem~\ref{thm:oscillatory-estimate}, applied
after this change of variables, bounds the sliced integral by
\[
 C_\nu\lambda^{-1+\nu}\|f(\cdot,y_3)\|_{L^3(\RR^2)}
\]
in $L^3$ of the output variables.  The transformed amplitudes have
uniform seminorms, and the output Jacobians are uniformly bounded.
Minkowski's inequality and then H\"older's inequality on the bounded
interval of $y_3$ prove the corresponding three-dimensional bound
for the selected cube pair.  Finally, if $a_k$ is the input $L^3$
norm in cube $k$, the output norm in cube $j$ is at most
$C_\nu\lambda^{-1+\nu}\sum_{k\sim j}a_k$.  The inequalities
\[
 \sum_j\left(\sum_{k\sim j}a_k\right)^3
 \le C\sum_k a_k^3
\]
follow from the bounded neighbor counts in both directions.  Summing
the finite angular partition proves \eqref{eq:br-unit-shell}.

\begin{proof}[Proof of Theorem~\ref{thm:bochner-riesz}]
Combining \eqref{eq:br-rescaled-shell} and \eqref{eq:br-unit-shell}
shows that the shell of radius $\lambda$ has $L^3$ operator norm at
most $C_{\delta,\nu}\lambda^{-\delta+\nu}$.  For each fixed
$\delta>0$, choose $\nu=\delta/2$.  The sum over $\lambda=2^j$
converges, and the bounded kernel piece is controlled by Young's
inequality.  The resulting operator is bounded on $L^3(\RR^3)$.

For a Schwartz input, convolution with $K_\delta$ equals
$\mathcal F^{-1}(m_\delta\widehat f)$: writing $K_\delta$ as the
inverse transform, Fubini is justified by
$\|m_\delta\|_1\|f\|_1<\infty$.  The spatial partition also
converges to this convolution pointwise by dominated convergence,
since $K_\delta$ is bounded and $f\in L^1$.  Its operator-norm sum
therefore represents the same multiplier.  Density of Schwartz
functions in $L^3$ supplies the asserted bounded extension.
\end{proof}

\subsection{The full strict range}

We record the consequence for other exponents.  Its proof uses the
positive-order theorem at arbitrarily small, fixed orders; no
uniformity as the order tends to zero is needed.

\begin{lemma}[Complex orders from a real order]\label{lem:br-complex-order}
For $\Re z>0$ and $R>0$, let $T_{z,R}$ have multiplier
$(1-|\xi|^2/R^2)_+^z$, and write $T_z=T_{z,1}$.  Suppose $r>0$
and $T_r$ is bounded on $L^q$ for $1\le q<\infty$.  If $a>r$ and
$N$ is an integer with $N>r+1$, then
\begin{equation}\label{eq:br-complex-bound}
 \|T_{a+i\tau}\|_{q\to q}
 \le \|T_r\|_{q\to q}
 \prod_{k=0}^{N-1}\left(1+\frac{\tau^2}{(a-r+k)^2}\right)^{1/2}.
\end{equation}
The same conclusion holds for $q=\infty$ if $K_r\in L^1$.
\end{lemma}

\begin{proof}
Conjugation by the $L^q$ isometry $f(x)\mapsto R^{3/q}f(Rx)$ shows
that $\|T_{r,R}\|_{q\to q}=\|T_r\|_{q\to q}$.  For $\Re z>r$,
the beta integral gives
\begin{equation}\label{eq:br-beta-mixture}
 T_z=\frac{\Gamma(z+1)}{\Gamma(r+1)\Gamma(z-r)}
 \int_0^1(1-u)^{z-r-1}u^rT_{r,\sqrt u}\,du.
\end{equation}
Indeed, at a frequency with $s=|\xi|^2<1$ the integral of the
multiplier on the right is
\[
 \int_s^1(1-u)^{z-r-1}(u-s)^r\,du
 =(1-s)^z B(r+1,z-r);
\]
both sides vanish if $s\ge1$.  For finite $q$, the integral in
\eqref{eq:br-beta-mixture} converges absolutely in $L^q$ on every
input: dilation is strongly continuous for positive radii and the
scalar majorant $(1-u)^{a-r-1}u^r$ is integrable.  The multiplier
identity first holds on Schwartz inputs and then extends by density.
For $q=\infty$, integrate instead the dilated kernels
$u^{3/2}K_r(\sqrt u\,\cdot)$ in $L^1$; their norms are constant,
and their resulting Fourier transform gives the same identity.

Taking norms in \eqref{eq:br-beta-mixture} yields
\begin{equation}\label{eq:br-gamma-bound}
 \|T_{a+i\tau}\|_{q\to q}
 \le \|T_r\|_{q\to q}
 \frac{\Gamma(a-r)}{\Gamma(a+1)}
 \left|\frac{\Gamma(a+1+i\tau)}{\Gamma(a-r+i\tau)}\right|.
\end{equation}
One can bound this quotient without asymptotics.  The gamma recurrence
and another beta identity give
\[
 \frac{\Gamma(z+1)}{\Gamma(z-r)}
 =\prod_{k=0}^{N-1}(z-r+k)
 \frac{B(z+1,N-r-1)}{\Gamma(N-r-1)}.
\]
Since $|B(a+1+i\tau,N-r-1)|\le B(a+1,N-r-1)$, substitution in
\eqref{eq:br-gamma-bound} and cancellation prove
\eqref{eq:br-complex-bound}.  In particular the boundary norms grow
at most polynomially in $|\tau|$.
\end{proof}

\begin{corollary}[Strict Bochner--Riesz range in three dimensions]
\label{cor:full-bochner-riesz}
Let $1\le p\le\infty$ and
\begin{equation}\label{eq:full-br-range}
 \delta>\max\left\{0,\,3\left|\frac1p-\frac12\right|-\frac12\right\}.
\end{equation}
Then the multiplier $(1-|\xi|^2)_+^\delta$ is bounded on
$L^p(\RR^3)$.  The same bounds hold uniformly for every positive
radius after dilation.  This assertion concerns the strict
inequality \eqref{eq:full-br-range}.
\end{corollary}

\begin{proof}
First let $3<p<\infty$.  Put
\[
 \theta=1-\frac3p\in(0,1),\qquad c=\delta-\theta>0,
 \qquad z(w)=c+w.
\]
On the boundary $\Re w=0$, Lemma~\ref{lem:br-complex-order} with
$r=c/2$ and Theorem~\ref{thm:bochner-riesz} give a polynomially
growing $L^3$ bound for $T_{c+i\tau}$.  On $\Re w=1$, apply the
same lemma with $r=1+c/2>1$.  Lemma~\ref{lem:br-kernel} supplies
an integrable real-order kernel and hence a polynomially growing
$L^\infty$ bound for $T_{1+c+i\tau}$.

This is an admissible analytic family.  For Schwartz functions $f,g$
its scalar pairing is the integral over the unit ball of
\[
 (1-|\xi|^2)^{c+w}\widehat f(\xi)\overline{\widehat g(\xi)};
\]
on compact subsets of $\Re w>-c$ all its derivatives are dominated
by integrable powers of $|\log(1-|\xi|^2)|$.  On the closed strip
$0\le\Re w\le1$, its $L^2$ norm is at most one.  Approximation by
Schwartz functions therefore extends weak analyticity in the open
strip and scalar continuity on the closed strip to all $L^2$ pairs,
including bounded simple functions of finite support.  Indeed, the
scalar pairings converge uniformly on the strip by the $L^2$ bound.
These are admissible tests for analytic interpolation.
To make the polynomial boundary growth harmless,
fix $\kappa>0$ and use the analytic family
\[
 e^{\kappa(w-\theta)^2}T_{c+w}.
\]
Its boundary norms are bounded because the scalar factor has modulus
$e^{\kappa((\Re w-\theta)^2-(\Im w)^2)}$.  The analytic
interpolation theorem \cite{Stein1956}, applied between $L^3$ and
$L^\infty$, gives an $L^p$ bound at $w=\theta$, since
\[
 \frac1p=\frac{1-\theta}{3},
 \qquad c+\theta=\delta.
\]
The Gaussian factor equals one there, so this is the required bound
for $T_\delta$.

For $2\le p\le3$, interpolate the fixed real operator $T_\delta$
between its $L^2$ bound from Plancherel and its $L^3$ bound from
Theorem~\ref{thm:bochner-riesz}; here \eqref{eq:full-br-range} is
exactly $\delta>0$.  Its real multiplier is self-adjoint on Schwartz
functions, so duality proves the corresponding bounds for
$1<p<2$.  The symmetry of \eqref{eq:full-br-range} under
$p\leftrightarrow p'$ gives precisely the asserted range.
Finally, for $p=1$ or $p=\infty$ the condition is $\delta>1$, and
Lemma~\ref{lem:br-kernel} and Young's inequality give the result
directly.  Dilation gives the assertion about radii.
\end{proof}

For completeness, the norm-convergence assertion in the introduction follows
from the same uniform radius bounds. If $1\le p<\infty$ and
$\widehat f\in C_c^\infty$, then $T_{\delta,R}f\to f$ in the
Schwartz topology: on the fixed Fourier support, every derivative of
$(1-|\xi|^2/R^2)^\delta$ converges to the corresponding derivative
of $1$. Inverse Fourier transforms of $C_c^\infty$ functions are dense
in $L^p$. The uniform bounds and a density argument therefore give
$T_{\delta,R}f\to f$ in $L^p$ for all $f$ in that space whenever
\eqref{eq:full-br-range} holds.
